\section{Objectifs}
\begin{itemize}
\item Relier produit scalaire, longueurs, angles et orthogonalité.
\item Choisir la formule de produit scalaire adaptée à la configuration.
\item Utiliser les identités de normes et la formule d’Al-Kashi.
\end{itemize}

\section{Prérequis}
Vecteurs, coordonnées, trigonométrie dans le triangle rectangle.

\section{Activité de découverte}
Comparer plusieurs façons de calculer le travail d’une force constante conduit à une quantité dépendant à la fois des normes et de l’angle : le produit scalaire.

\section{Vocabulaire}
Produit scalaire, norme, orthogonalité, projection orthogonale, base orthonormée.

\section{Cours}
Pour deux vecteurs non nuls formant un angle $\theta$ :
\[
\vec u\cdot\vec v=\|\vec u\|\,\|\vec v\|\cos\theta.
\]
Le produit scalaire est symétrique et bilinéaire. Deux vecteurs sont orthogonaux si et seulement si leur produit scalaire est nul.

Dans une base orthonormée, si $\vec u=(x,y)$ et $\vec v=(x',y')$ :
\[
\vec u\cdot\vec v=xx'+yy'.
\]

On utilise aussi :
\[
\|\vec u+\vec v\|^2=\|\vec u\|^2+2\vec u\cdot\vec v+\|\vec v\|^2.
\]
Dans un triangle, la formule d’Al-Kashi s’en déduit.

\section{Exemples corrigés}
Avec $\vec u=(2,-1)$ et $\vec v=(3,6)$, on obtient $6-6=0$ : les vecteurs sont orthogonaux.

\section{Méthode}
\begin{enumerate}
\item Repérer les données disponibles.
\item Choisir coordonnées, cosinus, projection ou identités de normes.
\item Calculer puis interpréter géométriquement.
\end{enumerate}

\section{Erreurs fréquentes}
\begin{itemize}
\item Confondre produit scalaire et produit des normes.
\item Employer la formule coordonnée dans une base non orthonormée.
\item Oublier l’interprétation géométrique du résultat.
\end{itemize}

\section{Synthèse}
Le produit scalaire est un outil métrique central pour démontrer une orthogonalité, calculer un angle ou une longueur et résoudre des problèmes de géométrie plane.
